% ============================================================================
% HOJA DE APROBACIÓN DE JURADOS
% ============================================================================

\newpage
\thispagestyle{empty}
\begin{center}
    \MakeUppercase{\fontsize{14pt}{1pt}\selectfont\textbf{Universidad Nacional del Altiplano}}
\end{center}

\begin{center}
    \MakeUppercase{\fontsize{12pt}{1pt}\selectfont\textbf{Puno}}
\end{center}

\begin{center}
    \MakeUppercase{\fontsize{12pt}{1pt}\selectfont{Facultad de Ingeniería Mecánica Eléctrica, Electrónica y Sistemas}}
\end{center}

\begin{center}
    \MakeUppercase{\fontsize{12pt}{1pt}\selectfont{Escuela Profesional de Ingeniería de Sistemas}}
\end{center}

\begin{center}
    \fontsize{11pt}{1.5pt}\selectfont{TESIS}
\end{center}

\begin{center}
    \MakeUppercase{\fontsize{12pt}{1.5pt}\selectfont\textbf{DETECCIÓN DE ATAQUES MINORITARIOS EN REDES IoT MEDIANTE UN MODELO HÍBRIDO BASADO EN STACKING, XGBOOST Y LIGHTGBM}}
\end{center}

\begin{center}
    \MakeUppercase{\fontsize{12pt}{1.5pt}\selectfont\textbf{PUNO 2026}}
\end{center}

\begin{center}
    \fontsize{12pt}{1pt}\selectfont{Presentado por:}
\end{center}

\begin{center}
    \fontsize{12pt}{1pt}\selectfont\textbf{Bach. Carmen Nieves Apaza Condori}
\end{center}

\begin{center}
    \fontsize{12pt}{1pt}\selectfont{Para optar el título profesional de:}
\end{center}

\begin{center}
    \MakeUppercase{\fontsize{12pt}{1pt}\selectfont\textbf{Ingeniero de Sistemas}}
\end{center}

\begin{flushleft}
    \fontsize{11pt}{1pt}\selectfont\textbf{APROBADA POR:}
\end{flushleft}

\begin{flushleft}
\begin{tabular}{p{4.5cm} p{10cm}}
    \textbf{PRESIDENTE:} & Pablo Cesar Tapia Catacora \\
                           & Dr. Presidente de Jurado \\[0.5cm]
    \textbf{PRIMER MIEMBRO:} & Carlos Boris Sosa Maydana \\
                              & Dr. Primer miembro de Jurado \\[0.5cm]
    \textbf{SEGUNDO MIEMBRO:} & Ubaldo Allca Mamani \\
                               & Dr. Segundo miembro de Jurado \\[0.5cm]
    \textbf{DIRECTOR / ASESOR:} & Hugo Yosef Gomez Quispe \\
                                 & Dr. Asesor de Tesis \\
\end{tabular}
\end{flushleft}

\begin{flushleft}
\textbf{Área:} \hspace{0.5cm} Inteligencia Artificial \\[0.3cm]
\textbf{Tema:} \hspace{0.5cm} Detección de ataques minoritarios en redes IoT mediante modelo híbrido basado en Stacking, XGBoost y LightGBM
\end{flushleft}

\vfill